% Feel free to use this preamble as is and add more to it! Define your own lstset and lstdefinelanguage.

% Please use correct typesetting for Latin abbreviations:
\newcommand{\ie}{\textit{i.e.},}
\newcommand{\eg}{\textit{e.g.},}

\usepackage{listings}
% Here you can define the general look of your listings, e.g., whether to have line numbers, how to display key words, comments, strings etc. There are a lot of other options. Check out our papers (https://db.cs.uni-tuebingen.de/publications/) to see what listings designs are possible. Get creative here!
\lstset{
  basicstyle=\ttfamily,
  basewidth=0.5em,
  numbers=left,
  numberstyle=\tiny\color{black!50},
  numbersep=5pt,
  keywordstyle=[1]\bfseries,
  keywordstyle=[2]\color{blue!70},
  % emphstyle=\color{code-pink},
  commentstyle=\color{black!40},
  stringstyle=\color{green}
}

% I define my own shell language instead of using the predefined bash language because it knows very little keywords. There are many more options you can define for your own language.
\lstdefinelanguage{sql}{
morestring=[d]{"},
morecomment=[l]{--},
keywords=[1]{WITH, RECURSIVE, SELECT, FROM, WHERE, AS},
keywords=[2]{MIN},
}

% Add a floating environment such that listings also float to the top or bottom of pages.
\usepackage{float}
\newfloat{lstfloat}{htbp}{lop}
\floatname{lstfloat}{Listing}
\def\lstfloatautorefname{Listing} % needed for hyperref/auroref

% Easy references for figures, listings, and tables. No need to provide the correct names/numbers. The package takes care of everything.
\usepackage{cleveref}
\crefname{section}{Section}{Sections}
\crefname{subsection}{Section}{Sections}
\crefname{subsubsection}{Section}{Sections}
\crefname{appendix}{Appendix}{Appendices}
\crefname{equation}{Equation}{Equations}
\crefname{figure}{Figure}{Figures}
\crefname{table}{Table}{Tables}
\crefname{subfigure}{Figure}{Figures}
\crefname{subtable}{Table}{Tables}
\crefname{lstfloat}{Listing}{Listings}

\usepackage{tikz}
\usetikzlibrary {shapes.geometric}